% concatenated from Sum_LMT_2026.tex
\documentclass[11pt]{article}

%\usepackage[notref,notcite]{showkeys}
%\usepackage{refcheck}

\usepackage[T1]{fontenc}
\usepackage{lmodern}

\usepackage{amsmath, amssymb, amsthm} 
\usepackage[margin=2.5cm]{geometry}
\usepackage{color, graphicx}
\usepackage[caption=false]{subfig}% Support for small, `sub' figures and tables

\usepackage{etoolbox}
\makeatletter
\patchcmd{\@maketitle}{\LARGE \@title}{\LARGE\bfseries\@title}{}{}

\renewcommand{\@seccntformat}[1]{\csname the#1\endcsname.\quad}
\makeatother

\definecolor{darkblue}{rgb}{0,0,.5}

\usepackage{hyperref}
\hypersetup{
	colorlinks=true,		% false: boxed links; true: colored links
	linkcolor=darkblue,		% color of internal links
	citecolor=red,		% color of links to bibliography
	urlcolor=darkblue 		% color of external links
}


\makeatletter
\def\th@plain{%
	\thm@notefont{}% same as heading font
	\itshape % body font
}
\def\th@definition{%
	\thm@notefont{}% same as heading font
	\normalfont % body font
}
\renewcommand{\qedsymbol}{\ensuremath{\blacksquare}}
\renewenvironment{proof}[1][\proofname]{\par
	\normalfont
	\topsep0\p@\@plus3\p@ \trivlist
	\item[\hskip\labelsep\itshape
	#1\@addpunct{.}]\ignorespaces
}{%
	\qed\endtrivlist
}
\makeatother


\newtheorem{theorem}{Theorem}[section]
\newtheorem{lemma}[theorem]{Lemma}
\newtheorem{corollary}[theorem]{Corollary}
\newtheorem{proposition}[theorem]{Proposition}
\newtheorem{fact}[theorem]{Fact}
\newtheorem{assumption}[theorem]{Assumption}
\theoremstyle{definition}
\newtheorem{definition}[theorem]{Definition}
\theoremstyle{definition}
\newtheorem{example}[theorem]{Example}
\theoremstyle{definition}
\newtheorem{remark}[theorem]{Remark}

\renewcommand\theenumi{\rm (\roman{enumi})}
\renewcommand\theenumii{(\rm \alph{enumii})}
\renewcommand{\labelenumi}{\rm (\roman{enumi})}
\renewcommand{\labelenumii}{\rm (\alph{enumii})}

\usepackage[shortlabels]{enumitem}


\parskip    4pt
\tolerance  3000

\allowdisplaybreaks


%-------------------------------------------------------------------------
\newcommand{\N}{\ensuremath{\mathbb N}}
\newcommand{\R}{\ensuremath{\mathbb R}}
\newcommand{\ball}{\ensuremath{\mathbb B}}

\newcommand{\argmin}{\ensuremath{\operatorname*{argmin}}}
\newcommand{\inte}{\ensuremath{\operatorname{int}}}
\newcommand{\reli}{\ensuremath{\operatorname{ri}}}
\newcommand{\ran}{\ensuremath{\operatorname{ran}}}
\newcommand{\zer}{\ensuremath{\operatorname{zer}}}
\newcommand{\dom}{\ensuremath{\operatorname{dom}}}
\newcommand{\intdom}{\ensuremath{\operatorname{int}\operatorname{dom}}\,}
\newcommand{\epi}{\ensuremath{\operatorname{epi}}}
\newcommand{\gra}{\ensuremath{\operatorname{gra}}}
\newcommand{\Fix}{\ensuremath{\operatorname{Fix}}}
\newcommand{\Id}{\ensuremath{\operatorname{Id}}}
\newcommand{\prox}{\ensuremath{\operatorname{Prox}}}
\newcommand{\sgn}{\ensuremath{\operatorname{Sign}}}


\newcommand{\tcblack}{\textcolor{black}}
\newcommand{\tcb}{\textcolor{blue}}
\newcommand{\tcr}{\textcolor{red}}
\newcommand{\todo}[1]{\tcr{{\bf TODO:}~#1}}

\def\beq{\begin{equation}}
\def\eeq{\end{equation}}
\def\baq{\begin{eqnarray}}
\def\eaq{\end{eqnarray}}
\def\baqn{\begin{eqnarray*}}
\def\eaqn{\end{eqnarray*}}


\if
{
\usepackage[pageref]{backref}
\renewcommand*{\backrefalt}[4]{%
\ifcase #1 %
(Not cited)%
\or
(Cited on p.~#2)%
\else
(Cited on pp.~#2)%
\fi
}
\renewcommand*{\backreftwosep}{, }
\renewcommand*{\backreflastsep}{, }
}
\fi

\def\tcb{ \textcolor{blue}}
\def\tcr{ \textcolor{red}}
 
\usepackage{float}

\begin{document}

\title{Splitting Algorithms Using  Pairs Monotonicity to Solve the Sum of Two Non-monotone Operators Problem}

\author{
Ba Khiet Le\thanks{Analytical and Algebraic Methods in Optimization Research Group, Faculty of Mathematics and Statistics, Ton
Duc Thang University, Ho Chi Minh City, Vietnam. E-mail: \texttt{lebakhiet@tdtu.edu.vn}}, 
~
Zakaria Mazgouri\thanks{Sidi Mohammed Ben Abdellah University, 
National School of Applied Sciences, Fez, Morocco. %ORCID 0000-0003-0131-4741
E-mail: \texttt{zakariam511@gmail.com}}, 
~and~ 
Michel Th\' era\thanks{Mathematics and Computer Science Department, University of Limoges, 123 Avenue Albert Thomas,
87060 Limoges CEDEX, France. E-mail: \texttt{michel.thera@unilim.fr}}
}



%\date{\today}

\maketitle


\begin{abstract}
\noindent {This paper is concerned with solving the problem of finding a zero of the sum of a single-valued operator and a set-valued operator, neither of which is assumed to be monotone, in a real Hilbert space. Under the framework of pair monotonicity, we propose a generalized forward--backward algorithm and a generalized Tseng  algorithm based on transformed and warped resolvents associated with an auxiliary linear operator. We establish weak, strong and linear convergence of the proposed algorithms under suitable assumptions involving pair strong monotonicity, pair cocoercivity, and pair Lipschitz continuity. Finally, numerical experiments are presented to validate the theoretical results and demonstrate the effectiveness of the proposed algorithms.
} 
\end{abstract}


\paragraph{Keywords:}
Pair monotonicity, 
nonconvex programming,  sum inclusion problem, forward-backward algorithm, Tseng-type algorithm,  transformed resolvent, warped resolvent.




\paragraph{Mathematics Subject Classification (MSC 2020):}
 49J52, 49J53.  


\section{Introduction}
%In this paper, we will  \tcr{focus}  on the %study
%  inclusion problem
%\begin{equation}\label{main-pb}
%0\in A(x)+B(x),
%\end{equation}
%where $\mathcal{H}$ is a real Hilbert space, $A:\mathcal{H}\to\mathcal{H}$ is single-valued, and $B:\mathcal{H}\rightrightarrows\mathcal{H}$ is set-valued. \tcr{ Bacause this formulation models so   many applications, } problem \eqref{main-pb}, when $A$ and $B$ are maximally monotone operators, has been extensively studied in the literature,   \tcr{with numerical and analytical techniques}; see, e.g., \cite{Apidopoulos,Attouch,Bauschke,Bauschke1,bp,bw,Chen,Davis,Gibali,Giselsson,Giselsson1,He,Hong,Lions,Moursi,Svaiter,Tseng}.
%\tcb{
%In this paper, we focus on the inclusion problem
%\begin{equation}\label{main-pb}
%0\in A(x)+B(x),
%\end{equation}
%where $\mathcal{H}$ is a real Hilbert space, $A:\mathcal{H}\to\mathcal{H}$ is single-valued, and $B:\mathcal{H}\rightrightarrows\mathcal{H}$ is set-valued. Because this formulation arises naturally in nonconvex optimization and encompasses a wide range of applications, problem \eqref{main-pb}, when $A$ and $B$ are maximally monotone operators, has been extensively studied in the literature using both analytical and numerical techniques; see, e.g., \cite{Apidopoulos,Attouch,Bauschke,Bauschke1,bp,bw,Chen,Davis,Gibali,Giselsson,Giselsson1,He,Hong,Lions,Moursi,Svaiter,Tseng}.
%}
%However, many problems lack monotonicity, and our objective in this contribution is therefore to study \eqref{main-pb} without imposing monotonicity assumptions on either operator. %Such 
%The absence of monotonicity, however, precludes the direct application of classical splitting methods. A useful framework for recovering an appropriate monotonicity structure is provided by the notion of \emph{pair monotonicity},  originally introduced in \cite{acl} and further developed in \cite{LMT,LDT}.
%Pair monotonicity was originally introduced in \cite{acl} for the numerical solution of a broad class of quasi-variational inequalities (QVIs) and inclusions of the form
%\begin{equation}\label{1oper}
%0\in\mathcal{A}(x),
%\end{equation}
%where $\mathcal{A}:\mathcal{H}\to\mathcal{H}$ is a nonmonotone operator. The central idea is to identify an auxiliary mapping $v:\mathcal{H}\to\mathcal{H}$ such that $\mathcal{A}$ and $v$ satisfy a generalized monotonicity relation, namely, that the pair $(\mathcal{A},v)$ is monotone. The mapping $v$ is selected in accordance with the structure and analytical properties of the operators, with the aim of restoring a suitable monotonicity property. Consequently, its construction is inherently problem-dependent.
%
%When $\mathcal{A}$ is single-valued and $v$ is invertible, a forward algorithm can be employed to solve \eqref{1oper}; see \cite{acl}.
%When $\mathcal{A}$  is set-valued, a generalized proximal point algorithm (GPPA) can be applied; see \cite{LDT}.
In this paper, we focus on the inclusion problem
\begin{equation}\label{main-pb}
0\in A(x)+B(x),
\end{equation}
where $\mathcal{H}$ is a real Hilbert space, $A:\mathcal{H}\to\mathcal{H}$ is single-valued, and $B:\mathcal{H}\rightrightarrows\mathcal{H}$ is set-valued. This fundamental formulation arises in a broad range of problems in optimization, variational analysis, equilibrium theory, and related areas. In particular, when $A$ and $B$ are maximally monotone, \eqref{main-pb} provides the underlying framework for numerous splitting and proximal methods, and has been extensively investigated from both analytical and computational perspectives; see, e.g., \cite{Apidopoulos,Attouch,Bauschke,Bauschke1,Chen,cp,cw,Davis,Gibali,Giselsson,Giselsson1,He,Hong,Lions,Moursi,Svaiter,Tseng}.

A substantially more challenging setting arises when monotonicity is absent. Many important models, particularly those originating in nonconvex optimization and equilibrium problems, lead naturally to inclusions in which neither operator is monotone. Our objective is to address precisely this setting by studying \eqref{main-pb} without imposing monotonicity assumptions on either $A$ or $B$. This departure from the classical framework is not merely technical: the loss of monotonicity fundamentally alters the geometry of the problem and, in general, rules out the direct use of standard splitting and proximal techniques. It therefore calls for new structural conditions that can recover, at least partially, the monotonicity mechanisms underlying the convergence of these methods.

A promising framework in this direction is provided by the notion of \emph{pair monotonicity}, originally introduced in \cite{acl} and subsequently developed in \cite{LDT,LMT}. Rather than requiring an operator to be monotone in its own right, pair monotonicity seeks to identify an auxiliary mapping that interacts with the given operator so as to restore a generalized monotonicity structure. This viewpoint is particularly well suited to nonmonotone inclusions, since it allows one to exploit a hidden monotonicity mechanism without imposing restrictive monotonicity assumptions on the original operator.

More precisely, pair monotonicity was used in \cite{acl} to study a broad class of quasi-variational inequalities (QVIs) and inclusions of the form
\begin{equation}\label{1oper}
0\in\mathcal{A}(x),
\end{equation}
where $\mathcal{A}:\mathcal{H}\rightrightarrows\mathcal{H}$ is not necessarily monotone. The key idea is to construct an auxiliary mapping $v:\mathcal{H}\to\mathcal{H}$ such that the pair $(\mathcal{A},v)$ satisfies an appropriate generalized monotonicity relation. In this framework, $v$ acts as a structural companion to $\mathcal{A}$ and is chosen to reflect the particular analytical properties of the problem. Thus, rather than modifying the original operator or imposing global monotonicity, one seeks to reveal a suitable monotonicity structure through the interaction between $\mathcal{A}$ and $v$.

This perspective leads naturally to algorithmic developments. When $\mathcal{A}$ is single-valued and $v$ is invertible, the pair-monotonicity framework yields a forward-type algorithm for solving \eqref{1oper}; see \cite{acl}. For set-valued operators, this approach was extended in \cite{LDT} through a generalized proximal point algorithm ($\mathbf{GPPA})$, thereby providing a mechanism for treating nonmonotone inclusions within a proximal framework. These results demonstrate that pair monotonicity can serve as a unifying structural principle for designing convergent algorithms beyond the classical monotone setting.





Problem \eqref{1oper} arises naturally in unconstrained optimization. Consider
$$
\min_{x\in\mathcal{H}} f(x),
$$
where $f:\mathcal{H}\to\mathbb{R}\cup\{+\infty\}$ is not necessarily convex. Under appropriate assumptions, a first-order necessary optimality condition can be expressed as
$$
0\in\partial f(x),
$$
where $\partial f$ denotes an appropriate generalized subdifferential, such as the Mordukhovich or Clarke subdifferential. A natural extension is the constrained optimization problem
$$
\min_{x\in C} f(x)
=\min_{x\in\mathcal{H}}\bigl(f(x)+I_C(x)\bigr),
$$
where $C\subset\mathcal{H}$ is a convex set and $I_C$ denotes its indicator function. Under suitable qualification conditions, a necessary optimality condition takes the form
\begin{equation}\label{constr}
0\in\partial f(x)+N_C(x),
\end{equation}
where $N_C$ denotes the normal cone mapping. Although $N_C$ is maximally monotone, the generalized subdifferential $\partial f$ is generally nonmonotone when $f$ is nonconvex. Nevertheless, the pair-monotonicity framework offers a way to recover a useful monotonicity structure: it may be possible to identify an auxiliary mapping $v$ such that both pairs $(\partial f,v)$ and $(N_C,v)$ are monotone. This observation provides a direct connection between pair monotonicity and constrained nonconvex optimization.

More generally, consider the composite minimization problem
$$
\min_{x\in\mathcal{H}}\bigl(f(x)+g(x)\bigr),
$$
where $f,g:\mathcal{H}\to\mathbb{R}\cup\{+\infty\}$. Under appropriate qualification conditions, its first-order optimality condition leads naturally to the inclusion \eqref{main-pb}, with $A=\partial f$ and $B=\partial g$. Thus, both constrained optimization and composite nonconvex optimization fit naturally within the framework of nonmonotone sum inclusions. This connection motivates our study of \eqref{constr} and, more generally, of \eqref{main-pb} through pair monotonicity.

Existing developments based on pair monotonicity have primarily focused on general nonmonotone inclusions and generalized proximal point methods; see \cite{LDT,LMT}. The present work takes a complementary step by explicitly exploiting the sum structure in \eqref{main-pb}. Rather than applying a proximal method to the inclusion as a whole, we develop splitting schemes that treat the single-valued operator $A$ by a forward step and the set-valued operator $B$ through a generalized resolvent. This distinction is particularly important for structured applications in which $A$ is explicitly computable, while $B$ possesses a resolvent that can be evaluated efficiently, as is the case for many normal cone and subdifferential operators. Our approach therefore extends the pair-monotonicity framework from generalized proximal methods to genuinely splitting-based algorithms for nonmonotone sum inclusions, while retaining the computational advantages associated with the separate treatment of $A$ and $B$.

Our first contribution is a generalized forward--backward algorithm based on a transformed resolvent of $B$ ($\mathbf{Algorithm~1}$). We establish linear convergence under two alternative pair-monotonicity conditions: either $(A,v)$ or $(B,v)$ is strongly monotone; see Theorems~\ref{thm1} and~\ref{thm2}, respectively. We further establish strong convergence under the assumption that {$(A,v)$ or $(v,A)$} is pair cocoercive  see Theorem~\ref{thm3}. These results provide convergence guarantees for a forward--backward scheme even though neither of the original operators is required to be monotone.

%When $A$ is not surjective, the generalized forward--backward framework may no longer be directly applicable. To overcome this limitation, we introduce a modified Tseng-type forward--backward--forward algorithm ($\mathbf{Algorithm~2}$). Under the range condition
%$$
%\operatorname{ran}A\subset\operatorname{ran}v,
%$$
%we establish convergence of the associated approximate solutions and moreover, under the additional assumption that $v^{-1}$ is Lipschitz continuous, we prove weak convergence of the generated sequence; see Theorem \ref{thm4}. This second algorithm therefore complements the forward--backward scheme by addressing a setting in which surjectivity of the single-valued operator is unavailable.
When strong monotonicity or cocoercivity is absent, the generalized forward–backward framework may no longer be directly applicable. To overcome this limitation, we introduce a generalized Tseng forward–backward–forward algorithm ($\mathbf{Algorithm 2}$) for computing approximate solutions. In particular, we allow $v$ to be nonsurjective, which provides greater flexibility in the choice of $v$ and thereby significantly broadens the range of potential applications.
Moreover, under the additional assumption that $v^{-1}$ is Lipschitz continuous, we prove weak convergence of the generated sequence; see Theorem \ref{thm4}. 

Taken together, these results demonstrate that pair monotonicity can be effectively combined with classical splitting principles to handle a substantially broader class of nonmonotone inclusions. In particular, the proposed framework separates the structural role of pair monotonicity from the computational roles of the two operators, thereby allowing convergence analysis and algorithm design to proceed without requiring either $A$ or $B$ to be monotone individually.

The remainder of the paper is organized as follows. Section \ref{sec2} reviews the notation and basic concepts related to pair monotonicity that will be used throughout the paper. In Section \ref{sec3}, we introduce the generalized forward--backward and generalized Tseng algorithms and establish their convergence properties. Section \ref{sec4} presents numerical experiments illustrating the theoretical results. Finally, Section \ref{sec5} concludes the paper with a summary of the main findings and directions for future research.






%Problem \eqref{1oper} arises naturally in unconstrained optimization,
%$$
%\min_{x\in\mathcal{H}} f(x),
%$$
%where $f:\mathcal{H}\to\mathbb{R}\cup \{+\infty\}$  is not necessarily convex. The corresponding first-order necessary optimality condition can be written as
%$$
%0\in\partial f(x),
%$$
%where $\partial f$ denotes an appropriate generalized subdifferential, such as the Mordukhovich or Clarke subdifferential. A natural extension is the constrained optimization problem
%$$
%\min_{x\in C}f(x)
%=\min_{x\in\mathcal{H}}{f(x)+I_C(x)},
%$$
%where $C\subset\mathcal{H}$ is a convex set and $I_C $ denotes its indicator function. Under appropriate qualification conditions, a necessary optimality condition takes the form
%\begin{equation}\label{constr}
%0\in\partial f(x)+N_C(x),
%\end{equation}
%where $N_C$ denotes the normal cone mapping. Although $N_C$ is maximally monotone, the subdifferential $\partial f$ is generally nonmonotone when $f$ is nonconvex. It may nevertheless be possible to identify an auxiliary mapping  $v$  such that both pairs $(N_C,v)$ and $(\partial f,v)$ are monotone. More generally, the minimization problem
%$$
%\min_{x\in\mathcal{H}}\bigl(f(x)+g(x)\bigr),
%$$
%where $f,g:\mathcal{H}\to\mathbb{R}\cup{+\infty}$, leads naturally to the inclusion \eqref{main-pb} with $A=\partial f$ and $B=\partial g$. This observation motivates the study of \eqref{constr} and, more generally, of the inclusion problem \eqref{main-pb} within the framework of pair monotonicity.
%
%Existing developments based on pair monotonicity have mainly focused on general nonmonotone inclusions and generalized proximal point methods; see \cite{LMT,LDT}. The present work addresses a complementary aspect by explicitly exploiting the decomposition in \eqref{main-pb}. Rather than applying a proximal method to the entire inclusion, we develop splitting procedures that treat the single-valued operator $A$ by a forward step and the set-valued operator  $B$  through a generalized resolvent. This approach is particularly relevant for structured problems in which $ A$ is explicitly computable while $B$ admits an efficiently computable resolvent or projection. Thus, the main contribution of this paper is to extend pair-monotonicity techniques from generalized proximal methods to splitting schemes for nonmonotone sum inclusions, while preserving the distinct roles of the two operators.
%
%Our first contribution is a generalized forward--backward algorithm based on the transformed resolvent of $B$ (Algorithm~1). We establish its linear convergence when either $(A,v)$ or $ (B,v))$ is strongly monotone (Theorems~\ref{thm1} and~\ref{thm2}, respectively). We also prove strong convergence under the assumptions that $(A,v)$ is pair cocoercive and $A^{-1}$ is single-valued and continuous (Theorem~\ref{thm3}).
%
%When $A$ is not surjective, we introduce a modified Tseng-type forward--backward--forward algorithm (Algorithm~2). Assuming that
%$$
%\operatorname{ran}A\subset\operatorname{ran}v,
%$$
%we establish convergence of the associated approximate solutions (Theorem~\ref{thm4}). Moreover, under the additional assumption that $v^{-1}$ is Lipschitz continuous, we prove weak convergence of the generated sequence (Theorem~\ref{thm5}).
%
%The remainder of the paper is organized as follows. Section~\ref{sec2} reviews the necessary notation and basic concepts related to pair monotonicity. In Section~\ref{sec3}, we present the proposed algorithms and establish their convergence properties. Numerical experiments illustrating the theoretical results are reported in Section~\ref{sec4}. Finally, concluding remarks and directions for future research are presented in Section~\ref{sec5}.
%%%%%%%%%%%%%%%%%%%%%%%%%%%%%%%%%%%%%%%%%%%%%%%%%%%%%%%%%%
%%%%%%%%%%%%%%%%%%%%%%%%%%%%%%%%%%%%%%%%%%%%%%%%%%%%%%%%%%%%%

\section{{Preliminaries }}\label{sec2}

{
Throughout this paper, \( \mathcal{H} \) denotes a real Hilbert space endowed with the inner product \( \langle \cdot, \cdot \rangle \) and the induced norm \( \|\cdot\| \).  For
a set-valued operator \( F : \mathcal{H} \rightrightarrows \mathcal{H} \), the \emph{domain}, \emph{range}, \emph{graph}, \emph{set of zeros}, and \emph{set of fixed points} of \( F \) are defined, respectively, by }
\begin{align*}
\dom F &= \{x \in \mathcal{H} : F(x) \neq \varnothing\}, \\
\ran F &= \bigcup_{x \in \mathcal{H}} F(x), \\
\gra F &= \{(x, y) : x \in \mathcal{H}, y \in F(x)\}, \\
\zer F &= \{x \in \mathcal{H} : 0 \in F(x)\}, \\
\Fix F &= \{x \in \mathcal{H} : x \in F(x)\}.
\end{align*}  

The \emph{inverse} of \( F \) is defined by \( F^{-1}(y) = \{x \in \mathcal{H} : y \in F(x)\} \). It is easy to see that \( \dom F = \ran F^{-1} \) and \( \ran F = \dom F^{-1} \).  

The set-valued mapping \( F \) is said to be \emph{monotone} if for all \( (x, x^*), (y, y^*) \in \gra F \),
\[
\langle x^* - y^*, x - y \rangle \geq 0,
\]
and \emph{\(\alpha\)-strongly monotone} if \( \alpha \in (0, +\infty) \) and for all \( (x, x^*), (y, y^*) \in \gra F \),
\[
\langle x^* - y^*, x - y \rangle \geq \alpha \|x - y\|^2.
\] 

	Next, we  recall the concept of \emph{pair monotonicity}, introduced in \cite{acl}, which provides a natural extension of the classical notion of monotonicity, which is a useful tool in the analysis of non-monotone inclusion problems (see \cite{LDT,LMT}). Let \( F_1, F_2 : \mathcal{H} \rightrightarrows \mathcal{H} \) be two set-valued mappings, we define 

\begin{itemize}
\item
The pair \( (F_1, F_2) \) is \emph{monotone} if for all \( x, y \in \mathcal{H} \),
$$
\langle F_1(x) - F_1(y), F_2(x) - F_2(y) \rangle \geq 0,
$$
i.e., for all \( x, y \in \mathcal{H} \), \( x_1^* \in F_1(x) \), \( y_1^* \in F_1(y) \), \( x_2^* \in F_2(x) \), \( y_2^* \in F_2(y) \),
$$
\langle x_1^* - y_1^*, x_2^* - y_2^* \rangle \geq 0.
$$
\item
The pair \( (F_1, F_2) \) is \emph{\(\alpha\)-strongly monotone} if \( \alpha \in (0, +\infty) \) and for all \( x, y \in \mathcal{H} \),
%\beq\label{strong}
$$
\langle F_1(x) - F_1(y), F_2(x) - F_2(y) \rangle \geq \alpha \|x - y\|^2,
$$
%\eeq
i.e., for all \( x, y \in \mathcal{H} \), \( x_1^* \in F_1(x) \), \( y_1^* \in F_1(y) \), \( x_2^* \in F_2(x) \), \( y_2^* \in F_2(y) \),
$$
\langle x_1^* - y_1^*, x_2^* - y_2^* \rangle \geq \alpha \|x - y\|^2.
$$
\end{itemize}  

\begin{remark}
While classical monotonicity compares an operator with the identity mapping ($F_2=Id$), pair monotonicity replaces the identity by a suitable auxiliary mapping, allowing certain non-monotone operators to inherit monotonicity-like properties.   	
\end{remark}  

\begin{definition}
The pair $(A,v)$ is called Lipschitz continuous if there exists $L>0$ such that 
$$
\Vert Ax-Ay\Vert\le L \Vert v(x)-v(y)\Vert.
$$
\end{definition}

\begin{definition}
The pair $(A,v)$ is called co-coercive  if there exists $\alpha>0$ such that 
$$
\langle Ax-Ay, v(x) - v(y) \rangle \geq \alpha \|Ax - Ay\|^2.
$$
\end{definition}

For a set-valued operator \( F : \mathcal{H} \rightrightarrows \mathcal{H} \) and a linear $v: \mathcal{H} \to \mathcal{H}$ with \( \dom v = \mathcal{H} \), we recall the following definitions. The \emph{transformed resolvent} of $F$ with respect to $v$ \cite{LDT} is defined by     
\[
T_{F}^v = v \circ (F + v)^{-1},
\] 
while the \emph{warped resolvent} of $F$ with kernel $v$ \cite{bc} is defined by
\[
J_F^v = (F + v)^{-1} \circ v,
\] 
provided that \( \ran v \subseteq \ran(F + v) \). These generalized resolvents will serve as the main tools in the design of the algorithms introduced in the next section. 


We   now  recall the notion of \emph{$R$-continuity} for a set-valued mapping
$\mathcal{A}:\mathcal{H}\rightrightarrows\mathcal{H}$ at a point
$\bar{x}$ in its domain; see \cite{Le,LMT,LT}. This property is useful in
the convergence analysis of the algorithms developed below.

\begin{definition}\label{rdef}
The set-valued mapping \noindent $\mathcal{A}:\mathcal{H} \rightrightarrows \mathcal{H}$ is called {\sc R-continuous} at ${\bar x}\in\dom\,\mathcal{A}$ if there exist a number $\sigma>0$ and a nondecreasing function $\rho: \mathbb{R}^+\to \mathbb{R}^+$ satisfying $\lim_{r\to 0^+}\rho(r)=\rho(0)=0$ such that we have the inclusion
\begin{equation}\label{rcon}
\mathcal{A}(x) \subset \mathcal{A}({\bar x}  )+\rho(\Vert x-{\bar x}   \Vert)\ball\;\mbox{ for all }\;x\in  \ball({\bar x} ,\sigma)
\end{equation}
with the {\sc continuity modulus function} $\rho$ and {\sc radius} $\sigma$. When $\sigma=\infty$, $\mathcal{A}$ is  called {\sc globally R-continuous} at ${\bar x}$. The inclusion in \eqref{rcon} means that for each $y\in \mathcal{A}(x)$, there exists ${\bar y}  \in  \mathcal{A}({\bar x} )$ satisfying the estimate 
$$
\Vert y-{\bar y}  \Vert\le \rho(\Vert x-{\bar x}   \Vert)
$$ 
for all $x\in \ball({\bar x},\sigma)$.
\end{definition}

%\begin{definition}\label{rdef}
%A set-valued mapping $\mathcal{A}:\mathcal{H}\rightrightarrows\mathcal{H}$
%is said to be \emph{$R$-continuous} at
%$\bar{x}\in\dom\mathcal{A}$ if there exist $\sigma>0$ and a nondecreasing
%function
%\[
%\rho:\mathbb{R}_+\to\mathbb{R}_+
%\]
%such that
%\[
%\rho(0)=0,
%\qquad
%\lim_{r\to0^+}\rho(r)=0,
%\]
%and
%\begin{equation}\label{rcon}
%\mathcal{A}(x)
%\subset
%\mathcal{A}(\bar{x})
%+\rho\bigl(\|x-\bar{x}\|\bigr)\mathbb{B}
%\qquad
%\text{for all }x\in\mathbb{B}(\bar{x},\sigma),
%\end{equation}
%where $\mathbb{B}$ denotes the closed unit ball of $\mathcal{H}$.
%The function $\rho$ is called a \emph{continuity modulus}, and $\sigma$ is
%called an \emph{$R$-continuity radius}. If $\sigma=+\infty$, then
%$\mathcal{A}$ is said to be \emph{globally $R$-continuous} at $\bar{x}$.
%
%Equivalently, \eqref{rcon}  means that, for every
%$x\in\mathbb{B}(\bar{x},\sigma)$ and every $y\in\mathcal{A}(x)$, there
%exists $\bar{y}\in\mathcal{A}(\bar{x})$ such that
%\[
%\|y-\bar{y}\|
%\leq
%\rho\bigl(\|x-\bar{x}\|\bigr).
%\]
%\end{definition}

\begin{definition}\label{closed graph}
Let $\mathcal{A}:\mathcal{H}\rightrightarrows\mathcal{H}$ be a set-valued
mapping. We say that:

\begin{itemize}
    \item the graph of $\mathcal{A}$ is  \emph{sequentially   closed }at a  point 
    $\bar{x}\in\dom\mathcal{A}$ if, for any sequences
    $x_k\to\bar{x}$ and $y_k\in\mathcal{A}(x_k)$ with $y_k\to y$, one has
    $y\in\mathcal{A}(\bar{x})$;

    \item the graph of $\mathcal{A}$ is \emph{sequentially weak-to-strong
    closed} if, for any sequences $x_k\rightharpoonup x$ and
    $y_k\in\mathcal{A}(x_k)$ with $y_k\to y$, one has
    $y\in\mathcal{A}(x)$.
\end{itemize}
\end{definition}

\begin{theorem}\label{compact}
\cite{LMT1,LT}
Suppose that $\mathcal{A}:\mathcal{H}\rightrightarrows\mathcal{H}$ has a
closed graph at $0$ and is locally compact around $0$, in the sense that
there exists $\sigma>0$ such that
\[
\mathcal{A}\bigl(\sigma\mathbb{B}\setminus\{0\}\bigr)
\]
is contained in a compact subset of $\mathcal{H}$. Then $\mathcal{A}$ is
$R$-continuous at $0$.
\end{theorem}

We finally recall the classical weak convergence lemma of Opial, which will be used to establish the weak convergence of the sequence generated by Algorithm~2.

\begin{lemma}[Opial's Lemma {\cite{Opial}}]\label{lem:Opial}
Let $S$ be a nonempty subset of $\mathcal{H}$, and let $(x_n)_{n\in\mathbb{N}}$ be a sequence in $\mathcal{H}$. Suppose that:
\begin{enumerate}[(i)]
    \item for every $x^*\in S$, the limit
    \[
    \lim_{n\to\infty}\|x_n-x^*\|
    \]
    exists;
     \item every sequential weak cluster point of $(x_n)_{n\in\mathbb{N}}$ belongs to $S$.
\end{enumerate}
Then $(x_n)_{n\in\mathbb{N}}$ converges weakly to some point $x^\infty\in S$.
\end{lemma}





%%%%%%%%%%%%%%%%%%%%%%%%%%%%%%%%%%%%%%%%%%%%%%%%%%%%%%%%%%%%%
%%%%%%%%%%%%%%%%%%%%%%%%%%%%%%%%%%%%%%%%%%%%%%%%%%%%%%%%%%%%%
\section{Main results}\label{sec3}

In this section, we study the convergence analysis of  algorithms that we propose to solve (\ref{main-pb}).   Throughout this section, we assume that the solution set
\[
S := \{x\in \mathcal{H} : 0\in A(x)+B(x)\},
\]
of \eqref{main-pb} is nonempty.

\medskip

\noindent\textbf{Assumption 1.} There exists a  mapping $v:\mathcal{H}\to\mathcal{H}$ such that the pairs $(A,v)$ and $(B,v)$ are monotone.

\medskip

\noindent\textbf{Assumption 2.} There exists $\gamma>0$ such that
\[
\operatorname{ran}(\gamma B+v)=\mathcal{H}.
\]

\medskip

Under Assumption 2, the transformed resolvent $T_{\gamma B}^{v}$ is well defined. Moreover, if the pair $(B,v)$ is monotone, then $T_{\gamma B}^{v}$ is firmly nonexpansive, i.e.,
\[
\|T_{\gamma B}^{v}(u)-T_{\gamma B}^{v}(w)\|^2 \le \langle T_{\gamma B}^{v}(u)-T_{\gamma B}^{v}(v),u-w\rangle,
\qquad \forall,u,w\in\mathcal{H}.
\]
Consequently, $T_{\gamma B}^{v}$ is  nonexpansive, i.e.,
\[
\|T_{\gamma B}^{v}(u)-T_{\gamma B}^{v}(w)\|^2 \le \|u-w\|,
\qquad \forall,u,w\in\mathcal{H}.
\]
This property is fundamental to our analysis, as it constitutes the key ingredient in the convergence analysis of Algorithm 1. We refer the reader to \cite{LDT} for the proof and further properties of transformed resolvents.  First  we investigate a generalized forward--backward scheme based on the transformed resolvent.

\medskip

\noindent\textbf{Algorithm 1.}
\[
x_{k+1}\in(\gamma B+v)^{-1}\bigl(v(x_k)-\gamma A(x_k)\bigr), \; k\ge 0,\; x_0\in \mathcal{H}.
\]

\medskip

The following theorems  establishes the strong convergence and  linear convergence for the sequence generated by Algorithm 1.


\begin{theorem}\label{thm1} 
Suppose that Assumptions 1 and 2 hold. If $(A,v)$ is $\alpha$-strongly monotone and $A$ is $L$-Lipschitz continuous, with
\[
0<\gamma<\frac{2\alpha}{L^2},
\]
then the sequence $(x_k)$ generated by Algorithm~1 converges to the unique solution $x^*$.

Moreover:
\begin{enumerate}
\item[(a)] If $v$ is Lipschitz continuous, then $(v(x_k))$ converges linearly to $v(x^*)$.
\item[(b)] If both $v$ and $v^{-1}$ are Lipschitz continuous, then $(x_k)$ converges linearly to $x^*$.
\end{enumerate}
\end{theorem}

\begin{proof}
Since $(A,v)$ is $\alpha$-strongly monotone, the solution set $S$ has a unique point $x^*$.
{Since $v(x_{k+1})=T_{\gamma B}^v(v(x_k)-\gamma Ax_k)$ and $T_{\gamma B}^v$ is nonexpansive, we have
\baqn
\Vert v(x_{k+1})-v(x^*)\Vert^2&\le&\Vert v(x_{k})-v(x^*)-\gamma (Ax_k-Ax^*)\Vert^2\\
&=&\Vert v(x_{k})-v(x^*)\Vert^2-2\gamma\langle v(x_{k})-v(x^*)), Ax_k-Ax^*\rangle+\gamma^2 \Vert Ax_k-Ax^*\Vert^2\\
&\le&\Vert v(x_{k})-v(x^*)\Vert^2-2\gamma \alpha \Vert x_{k}-x^*\Vert^2+\gamma^2L^2\Vert x_{k}-x^*\Vert^2\\
&\le&\Vert v(x_{k})-v(x^*)\Vert^2-\beta \Vert x_{k}-x^*\Vert^2,
\eaqn
where
\[
\beta:=2\gamma\alpha-\gamma^2L^2>0.
\]
Hence, the sequence $\bigl(\|v(x_k)-v(x^*)\|\bigr)$ is monotonically decreasing and therefore convergent.  Consequently $$\|x_k-x^*\|\to0.$$
}
%\medskip

\noindent
 {
(a) Suppose that $v$ is Lipschitz continuous with modulus $L_v$. Then
\[
\|v(x_k)-v(x^*)\|
\le
L_v\|x_k-x^*\|.
\]
%which implies
%$$
%\Vert v(x_{k+1})-v(x^*)\Vert^2\le\Vert v(x_{k})-v(x^*)\Vert^2-\frac{\beta}{L_v^2} \Vert v(x_{k})-v(x^*)\Vert^2=(1-\frac{\beta}{L_v^2})\Vert v(x_{k})-v(x^*)\Vert^2.
%$$
Substituting this estimate into the previous inequality gives
\[
\|v(x_{k+1})-v(x^*)\|^2
\le
\left(1-\frac{\beta}{L_v^2}\right)
\|v(x_k)-v(x^*)\|^2.
\]
Therefore, $(v(x_k))$ converges linearly to $v(x^*)$.
}
\medskip

\noindent
{(b) Assume, in addition, that $v^{-1}$ is Lipschitz continuous with modulus $\widetilde L_v$. Then
\[
\|x_{k+1}-x^*\|
\le
\widetilde L_v\|v(x_{k+1})-v(x^*)\|
\le
\widetilde L_v
\sqrt{1-\frac{\beta}{L_v^2}}
\|v(x_k)-v(x^*)\|
\le
L_v\widetilde L_v
\sqrt{1-\frac{\beta}{L_v^2}}
\|x_k-x^*\|.
\]
Hence, $(x_k)$ converges linearly to $x^*$.
}
\end{proof}

%\newpage 

\begin{theorem}\label{thm2}
Suppose that Assumptions 1 and 2 hold. Assume that the pair $(B,v)$ is $\alpha$-strongly monotone,   the pair $(A,v)$ is $L$-Lipschitz continuous and $v$ is $L_v$-Lipschitz continuous for some $L, L_v>0$. Then the sequence $(v(x_k))$ generated by Algorithm 1 converges linearly to $v(x^*)$, provided that
\[
\gamma<\alpha L^{-2}L_v^{-2}.
\]
\end{theorem}

\begin{proof}
Since
\[
v(x_{k+1})=
T_{\gamma B}^{v}\bigl(v(x_k)-\gamma A(x_k)\bigr),
\]
and $T_{\gamma B}^{v}$ is
$\frac{1}{1+\alpha\gamma L_v^{-2}}$-Lipschitz continuous (see \cite{LDT}), it follows that
\baqn
\Vert v(x_{k+1})-v(x^*)\Vert^2&\le&\frac{1}{1+\alpha \gamma L_v^{-2}} \Vert v(x_{k})-v(x^*)-\gamma (Ax_k-Ax^*)\Vert^2\\
&=&\frac{1}{1+\alpha \gamma L_v^{-2}}(\Vert v(x_{k})-v(x^*)\Vert^2-2\gamma\langle v(x_{k})-v(x^*), Ax_k-Ax^*\rangle\\
&&\,+\gamma^2 \Vert Ax_k-Ax^*\Vert^2).
%&\le&\frac{1}{1+\alpha \gamma L^{-2}}(1+\gamma^2L^2)\Vert v(x_{k})-v(x^*)\Vert^2\\
%&\le&\kappa \Vert v(x_{k})-v(x^*)\Vert^2.
\eaqn   
Since $(A,v)$ is $L$-Lipschitz continuous, one has 
\[
\|Ax_k-Ax^*\|
\le
L\|v(x_k)-v(x^*)\|.
\]
Discarding the nonpositive cross term yields
\[
\|v(x_{k+1})-v(x^*)\|^2
\le
\frac{1+\gamma^2L^2}
{1+\alpha\gamma L_v^{-2}}
\|v(x_k)-v(x^*)\|^2=\kappa \|v(x_k)-v(x^*)\|^2
\]
where
\[
\kappa:=
\frac{1+\gamma^2L^2}
{1+\alpha\gamma L_v^{-2}}<1
\]
since
\[
\gamma<\alpha L^{-2}L_v^{-2}.
\]
The conclusion follows.
\end{proof}
% 

%
%\begin{theorem}\label{thm3} 
%Suppose that Assumptions 1, 2 are satisfied. \\
%a) If $(A,v)$ is $\alpha$-cocoercive and  $\gamma<2\alpha$ then the sequence $(A(x_k))$ generated by Algorithm 1 converges to $A(x^*)$. In addition, if $A^{-1}$ is single-valued continuous then $(x_k)$ converges to $x^*$.\\
%\tcb{b) If $(v,A)$ is $\alpha$-cocoercive, $(A,v)$ is $L$-Lipschitz continuous and  $\gamma<2\alpha/L^2$ then the sequence $(v(x_k))$ generated by Algorithm 1 converges linearly to $v(x^*)$. In addition, if $v^{-1}$ is single-valued Lipschitz continuous then $(x_k)$ converges linearly to $x^*$.}
%\end{theorem}

\begin{theorem}\label{thm3}
Suppose that Assumptions~1 and~2 are satisfied.

\begin{enumerate}[(a)]
\item If $(A,v)$ is $\alpha$-cocoercive and $\gamma<2\alpha$, then the
sequence $(A(x_k))$ generated by Algorithm~1 converges to $A(x^*)$.
Moreover, if $A^{-1}$ is single-valued and continuous, then $(x_k)$
converges to $x^*$.

\item If $(v,A)$ is $\alpha$-cocoercive, $(A,v)$ is
$L$-Lipschitz continuous, and $\gamma<2\alpha/L^2$, then the sequence
$(v(x_k))$ generated by Algorithm~1 converges linearly to $v(x^*)$.
Moreover, if $v^{-1}$ is single-valued and Lipschitz continuous, then
$(x_k)$ converges linearly to $x^*$.
\end{enumerate}
\end{theorem}


\begin{proof}
a) Since $T_{\gamma B}^v$ is nonexpansive, one has 
\baqn
\Vert v(x_{k+1})-v(x^*)\Vert^2&\le&\Vert v(x_{k})-v(x^*)-\gamma (Ax_k-Ax^*)\Vert^2\\
&=&\Vert v(x_{k})-v(x^*)\Vert^2-2\gamma\langle v(x_{k})-v(x^*)), Ax_k-Ax^*\rangle+\gamma^2 \Vert Ax_k-Ax^*\Vert^2\\
&\le&\Vert v(x_{k})-v(x^*)\Vert^2-2\gamma \alpha \Vert Ax_k-Ax^*\Vert^2+\gamma^2\Vert Ax_k-Ax^*\Vert^2\\
&\le&\Vert v(x_{k})-v(x^*)\Vert^2-\beta \Vert Ax_k-Ax^*\Vert^2,
\eaqn
where $\beta=\gamma(2\alpha-\gamma)>0$.
Therefore $(\Vert v(x_{k+1})-v(x^*)\Vert)$ is decreasing, convergent and  $(\Vert Ax_k-Ax^*\Vert)$ converges to $0$. The continuity of  $A^{-1}$ implies the convergence of the sequence $(x_k)$ to $x^*$.\\
%\tcb{b) Similarly, we have
%\baqn
%\Vert v(x_{k+1})-v(x^*)\Vert^2&\le&\Vert v(x_{k})-v(x^*)\Vert^2-2\gamma\langle v(x_{k})-v(x^*)), Ax_k-Ax^*\rangle+\gamma^2 \Vert Ax_k-Ax^*\Vert^2\\
%&\le&\Vert v(x_{k})-v(x^*)\Vert^2-2\gamma \alpha \Vert v(x_k)-v(x^*)\Vert^2+\gamma^2L^2\Vert v(x_k)-v(x^*)\Vert^2\\
%&\le&(1-\kappa)\Vert v(x_{k})-v(x^*)\Vert^2,
%\eaqn
%where $\kappa=\gamma(2\alpha-\gamma L^2)>0$. The conclusion follows.}

(b) Similarly, using the nonexpansiveness of $T_{\gamma B}^v$, we obtain
\begin{align*}
\|v(x_{k+1})-v(x^*)\|^2
&\leq
\|v(x_k)-v(x^*)\|^2
-2\gamma\langle v(x_k)-v(x^*),Ax_k-Ax^*\rangle
+\gamma^2\|Ax_k-Ax^*)\|^2\\
&\leq
\|v(x_k)-v(x^*)\|^2
-2\gamma\alpha\|v(x_k)-v(x^*)\|^2
+\gamma^2L^2\|v(x_k)-v(x^*)\|^2\\
&=
\bigl(1-\gamma(2\alpha-\gamma L^2)\bigr)
\|v(x_k)-v(x^*)\|^2\\
&=(1-\kappa)\|v(x_k)-v(x^*)\|^2,
\end{align*}
where
\[
\kappa:=\gamma(2\alpha-\gamma L^2)>0.
\]
Since $\gamma<2\alpha/L^2$, we also have $0<\kappa<1$. Hence,
\[
\|v(x_k)-v(x^*)\|
\leq
(1-\kappa)^{k/2}\|v(x_0)-v(x^*)\|,
\]
which proves the linear convergence of $(v(x_k))$ to $v(x^*)$.
If, in addition, $v^{-1}$ is single-valued and Lipschitz continuous with
Lipschitz modulus $\widetilde L_v$, then
\[
\|x_k-x^*\|
\leq
\widetilde L_v\|v(x_k)-v(x^*)\|
\leq
\widetilde L_v(1-\kappa)^{k/2}
\|v(x_0)-v(x^*)\|,
\]
and therefore $(x_k)$ converges linearly to $x^*$.
\end{proof}

Algorithm 1 relies on strong monotonicity or cocoercivity and on the full range of $v$, which may be restrictive in certain applications. The following weaker assumption allows us to construct a generalized Tseng forward–backward–forward algorithm that remains applicable under weaker conditions and provides approximate solutions to the inclusion problem.

%
%\tcb{ Algorithm 1 require the global range condition of Assumption 2 and the strong monotonicity or co-coercivity, which may be restrictive in certain applications. The following weaker assumption enables the construction of a Tseng-type forward-backward algorithm to find approximate solutions of the inclusion problem.} \\
%
\noindent \textbf{Assumption 3.} The mapping $v:\mathcal{H}\to\mathcal{H}$ is linear % \tcr{ and injective } 
and $\ran  A\subset  \ran v\subset \ran (\gamma B+v)$. \\

{Under Assumption 3, the mappings $v^{-1}A$ and $J_{\gamma B}^v$ are well defined. This allows us to introduce the following Tseng-type forward-backward algorithm.}\\

\noindent {\bf Algorithm 2}.
\[
\begin{cases}
y_{k}\in (\gamma B+v)^{-1}(v(x_k)-\gamma Ax_k)=J^v_{\gamma B}(x_k-\gamma v^{-1}(Ax_k)),\\
x_{k+1}\in y_k+\gamma v^{-1}(Ax_k-Ay_k), \;k\ge 0, \;x_0\in \mathcal{H}.
\end{cases}
\]

\medskip
\begin{theorem}\label{thm4}
Suppose that Assumptions 1 and 3 hold, and that the pair $(A,v)$ is $L$-Lipschitz continuous for some $L>0$ satisfying
$$
\gamma<\frac{1}{L}.
$$
Let $(x_k)$ and $(y_k)$ be the sequences generated by Algorithm 2. Then the following assertions hold:
\begin{enumerate}
\item
$
\|v(y_k)-v(x_k)\|\to0\qquad\text{as }k\to\infty.
$

\item If $y^*$ is a weak cluster point of $(y_k)$ and the graph of $A+B$ is sequentially weak-to-strong closed, then $y^*\in S$. Moreover, if $v^{-1}:\operatorname{ran}v\to\mathcal{H}$ is Lipschitz continuous, then the sequence $(x_k)$ generated by Algorithm 2 converges weakly to a point in $S$.

\item If $(A+B)^{-1}$ is $R$-continuous at $0$, then
$
d(y_k,S)\to0.
$
\end{enumerate}
\end{theorem}




%
%
%\begin{theorem}\label{thm4}
%Suppose that Assumptions~1 and~3 hold, and that the pair $(A,v)$ is $L$-Lipschitz continuous for some $L>0$ satisfying
%\[
%\gamma<\frac{1}{L}.
%\]
%Let $(x_k)$ and $(y_k)$ be the sequences generated by Algorithm~2. Then the following assertions hold:
%\begin{enumerate}
%\item$\|v(y_k)-v(x_k)\|\to0$ as $k\to\infty$.
%
%\item If $y^*$ is a weak cluster point of $(y_k)$ and the graph of $(A+B) $ is sequentially weak-to-strong closed, then $y^*\in S$.
%
%\item If $(A+B)^{-1}$ is $R$-continuous at $0$, then
%$
%d(y_k,S)\to0.
%$
%\end{enumerate}
%\end{theorem}
\begin{proof}
(i) Let $x^*\in S$. Then
\begin{equation}\label{discr}
\frac{v(y_k)-v(x_k)}{\gamma}+A(x_k)\in -B(y_k),
\end{equation}
and
\[
A(x^*)\in -B(x^*).
\]
Since $(B,v)$ is monotone, it follows that
\[
\left\langle
\frac{v(y_k)-v(x_k)}{\gamma}+A(x_k)-A(x^*),
v(y_k)-v(x^*)
\right\rangle
\le0.
\]
Combining this inequality with the monotonicity of $(A,v)$, we obtain
\begin{equation}\label{eq0}
\left\langle
\frac{v(y_k)-v(x_k)}{\gamma}+A(x_k)-A(y_k)
,v(y_k)-v(x^*)
\right\rangle
\le0.
\end{equation}
On the other hand, Algorithm~2 yields
\begin{equation}\label{eq00}
A(x_k)-A(y_k)=
\frac{1}{\gamma}\bigl(v(x_{k+1})-v(y_k)\bigr).
\end{equation}
Substituting \eqref{eq00} into \eqref{eq0} and using $\gamma>0$, we obtain
\[
\langle
v(x_{k+1})-v(x_k)
,v(y_k)-v(x^*)
\rangle
\le0.
\]
Hence,
\baqn
	&&\langle v(x_{k+1})-v(x^*), v(y_k)-v(x^*)\rangle\le \langle v(x_{k})-v(x^*),\, v(y_k)-v(x^*)\rangle\\
	&\Leftrightarrow& \Vert v(x_{k+1})-v(x^*)\Vert^2+\Vert v(y_k)-v(x^*)\Vert^2-\Vert v(x_{k+1})-v(y_k)\Vert^2\\
	&\le&\Vert v(x_{k})-v(x^*)\Vert^2+\Vert v(y_k)-v(x^*)\Vert^2-\Vert v(x_{k})-v(y_k)\Vert^2\\
	&\Leftrightarrow& \Vert v(x_{k+1})-v(x^*)\Vert^2\le \Vert v(x_{k})-v(x^*)\Vert^2+\gamma^2 \Vert Ax_k-Ay_k\Vert^2-\Vert v(x_{k})-v(y_k)\Vert^2,
	\eaqn
where the last equality follows from \eqref{eq00}.

Since $(A,v)$ is $L$-Lipschitz continuous and
$\gamma^2L^2<1,
$ there exists $\varepsilon>0$ such that
$
\gamma^2L^2\le1-\varepsilon.
$
Therefore,
\[
\|v(x_{k+1})-v(x^*)\|^2
\le
\|v(x_k)-v(x^*)\|^2
-\varepsilon\|v(x_k)-v(y_k)\|^2.
\]
It follows that the sequence
$(\Vert v(x_{k+1})-v(x^*)\Vert)$  is nonincreasing and hence convergent.  Moreover,
\[
\|v(x_k)-v(y_k)\|\longrightarrow0.
\]

\medskip

\noindent
(ii) Let $y^*$ be a weak cluster point of $(y_k)$. By \eqref{discr},
\begin{equation}\label{discr1}
\frac{v(y_k)-v(x_k)}{\gamma}
+A(x_k)-A(y_k)
\in
-(A+B)(y_k).
\end{equation}
Since
\[
\|v(x_k)-v(y_k)\|\to0,
\]
and $(A,v)$ is $L$-Lipschitz continuous, it follows that
\[
\frac{v(y_k)-v(x_k)}{\gamma}
+A(x_k)-A(y_k)
\longrightarrow0.
\]
Using the assumption that the graph of $A+B$ is sequentially weak-to-strong closed, we conclude that
\begin{equation} \label{belong}
0\in A(y^*)+B(y^*),
\end{equation}
that is,
$
y^*\in S.
$

\medskip
%Now suppose that $v^{-1}:\operatorname{ran}v\to\mathcal{H}$ is Lipschitz continuous. Since the sequence
%$$
%\bigl(\|v(x_{k+1})-v(x^*)\|\bigr)
%$$
%converges and
%$
%\|v(x_k)-v(y_k)\|\to0,
%$
%the Lipschitz continuity of $v^{-1}$ implies that
%$$
%\bigl(\|x_{k+1}-x^*\|\bigr)
%$$
%also converges and that
%$
%\|x_k-y_k\| \to 0.
%$
%Let $x^*$ be a weak cluster point of $(x_k)$. Since $\|x_k-y_k\|\to0$, $x^*$ is also a weak cluster point of $(y_k)$. By \eqref{belong}, we then have $x^*\in S$. Finally, Opial's lemma yields that $(x_k)$ converges weakly to some point in $S$.
Now suppose that $v^{-1}:\operatorname{ran}v\to\mathcal{H}$ is Lipschitz continuous. Since the sequence
$$
\bigl(\|v(x_{k+1})-v(x^*)\|\bigr)
$$
converges and
$$
\|v(x_k)-v(y_k)\|\to0,
$$
the Lipschitz continuity of $v^{-1}$ implies that
$
\bigl(\|x_{k+1}-x^*\|\bigr)
$
also converges and that
$
\|x_k-y_k\| \to 0.
$
Let $x^*$ be a weak cluster point of $(x_k)$. Since $\|x_k-y_k\|\to0$, $x^*$ is also a weak cluster point of $(y_k)$. By \eqref{belong}, it follows that $x^*\in S$. Finally, Opial's lemma implies that $(x_k)$ converges weakly to some point in $S$.

\noindent
(iii) Define
\[
u_k:=
\frac{v(y_k)-v(x_k)}{\gamma}
+A(x_k)-A(y_k).
\]
Then $u_k\to0$. By \eqref{discr1},
\[
y_k\in(A+B)^{-1}(-u_k).
\]
Since $(A+B)^{-1}$ is $R$-continuous at $0$, we have
\[
(A+B)^{-1}(-u_k)
\subset
(A+B)^{-1}(0)+\rho(\|u_k\|),
\]
where $\rho(t)\to0$ as $t\to0$. Consequently,
\[
y_k
\in
S+\rho(\|u_k\|),
\]
which implies that
\[
d(y_k,S)\longrightarrow0.
\]
This completes the proof.
\end{proof}

\begin{remark}
	\label{rem:3.5}
	\rm(a) At the end of (ii) in Theorem \ref{thm4}, if
	$v^{-1}:\operatorname{ran}v\to \mathcal{H}$ is not single-valued and Lipschitz continuous, {but there exists a single-valued Lipschitz continuous selection $\widetilde v^{-1}$} , then $v^{-1}$ can be replaced by $\widetilde v^{-1}$ and the same conclusion follows.
	
	\rm(b) If the range condition
	\[
	\operatorname{ran} A\subset \operatorname{ran}v
	\subset \operatorname{ran}(\gamma B+v)
	\]
	is not satisfied, one may consider an alternative decomposition
	\[
	A+B=A'+B'
	\]
	such that, for a suitable choice of $v$, the corresponding range condition
	\[
	\operatorname{ran} A'\subset \operatorname{ran}v
	\subset \operatorname{ran}(\gamma B'+v)
	\]
	holds. This provides additional flexibility in applying the generalized framework.
\end{remark}

%
%\tcr{
%\begin{theorem}\label{thm5}
%Suppose that the assumptions of Theorem~\ref{thm4} hold and that
%$v^{-1}:\operatorname{ran}v\to\mathcal{H}$ is Lipschitz continuous.
%Then the sequence $(x_k)$ generated by Algorithm~2 converges weakly to
%some $x^*\in S$.
%\end{theorem}
%}
%
%\begin{proof}
%By Theorem~\ref{thm4}(i),
%\[
%\|v(x_k)-v(y_k)\|\longrightarrow 0.
%\]
%Since $v^{-1}:\operatorname{ran}v\to\mathcal{H}$ is Lipschitz continuous
%with modulus $L_{v^{-1}}$, we have
%\[
%\begin{aligned}
%\|x_k-y_k\|
%&=
%\|v^{-1}(v(x_k))-v^{-1}(v(y_k))\|\\
%&\leq
%L_{v^{-1}}\|v(x_k)-v(y_k)\|
%\longrightarrow0.
%\end{aligned}
%\]
%Hence, $(x_k)$ and $(y_k)$ have the same weak cluster points.
%
%Set
%\[
%z_k:=v(x_k).
%\]
%By \eqref{discr0}, for every $x^*\in S$ there exists $\varepsilon>0$
%such that
%\[
%\|z_{k+1}-v(x^*)\|^2
%\leq
%\|z_k-v(x^*)\|^2
%-\varepsilon\|v(x_k)-v(y_k)\|^2.
%\]
%Consequently, for every $x^*\in S$, the sequence
%\[
%\bigl(\|z_k-v(x^*)\|\bigr)_{k\in\mathbb N}
%\]
%is nonincreasing and therefore convergent. In particular, $(z_k)$ is
%bounded.
%
%Let $z$ be a weak cluster point of $(z_k)$. Then, along some subsequence
%$(z_{k_j})$,
%\[
%z_{k_j}\rightharpoonup z.
%\]
%Since $v$ is linear and $v^{-1}:\operatorname{ran}v\to\mathcal{H}$ is
%Lipschitz continuous, $v^{-1}$ is linear and continuous on
%$\operatorname{ran}v$. Moreover, the Lipschitz continuity of $v^{-1}$
%implies that $\operatorname{ran}v$ is closed. Hence
%$z\in\operatorname{ran}v$, and
%\[
%x_{k_j}
%=
%v^{-1}(z_{k_j})
%\rightharpoonup
%v^{-1}(z).
%\]
%Since
%\[
%\|x_{k_j}-y_{k_j}\|\longrightarrow0,
%\]
%we also have
%\[
%y_{k_j}\rightharpoonup v^{-1}(z).
%\]
%By Theorem~\ref{thm4}(ii), every weak cluster point of $(y_k)$ belongs
%to $S$. Therefore,
%\[
%v^{-1}(z)\in S,
%\]
%and hence
%\[
%z\in v(S).
%\]
%Thus, every weak cluster point of $(z_k)$ belongs to $v(S)$.
%
%On the other hand, for every $z^*\in v(S)$, there exists $x^*\in S$
%such that $z^*=v(x^*)$. By the preceding argument,
%\[
%\|z_k-z^*\|
%=
%\|v(x_k)-v(x^*)\|
%\]
%has a finite limit as $k\to\infty$. Therefore, Opial's lemma implies
%that $(z_k)$ converges weakly to some $z^*\in v(S)$.
%
%Finally, since $v^{-1}$ is continuous and linear on $\operatorname{ran}v$,
%we obtain
%\[
%x_k=v^{-1}(z_k)
%\rightharpoonup
%v^{-1}(z^*)
%=:x^*.
%\]
%Since $z^*\in v(S)$, we have $x^*\in S$. This proves the result.
%\end{proof}
%
%

%\section{Numerical Examples}\label{sec4}
%
%\begin{example}
%First we consider the operators $A: \R^3\to \R^3, B: \R^3 \rightrightarrows \R^3$ defined by
%$$
%A(x)=Qx+b, \qquad B(x)= N_{\R^2_+\times \R_-}(x),
%$$
%where
%$$
%Q =\begin{bmatrix} 5 & -1 & -1 \\ 1 & 4 & 0 \\ 1 & 1 & -4 \end{bmatrix},  \quad
%b = \begin{bmatrix} -1 \\ -1 \\ 1 \end{bmatrix}.
%$$
%\medskip
%We choose 
%$$
%v=\begin{bmatrix} 2 & 0 & 0\\ 0 & 3 & 0 \\ 0 & 0 & -2\end{bmatrix}.
%$$
%Then $A$ is non-monotone,  $v$ is invertible, $(B,v)$ is monotone and $(A,v)$  is strongly monotone. Moreover, both $v$ and $v^{-1}$ are Lipschitz continuous. Thus, all the assumptions of Theorem~\ref{thm1}
%are satisfied, and the corresponding convergence result applies. A direct computation yields the unique solution $x^*= \tcb{???}$.
%
%
%\medskip
%
%\noindent
%
%
%To illustrate the convergence behavior of Algorithm 1, we compute the errors $\|x_k-x^*\|_2$ and $\|v(x_k)-v(x^*)\|_2$. The numerical result in Figure \ref{error1} illustrates both errors when $\gamma=0.3$ and $x_0=(2,2,2)$.
%
%\begin{figure}[h!]
%	\centering{\includegraphics[scale=0.4]{fixed gamma and x0.png}}
%	\caption{\centering {\small The iterate errors $\|x_k-x^*\|_2$ and $\|v(x_k)-v(x^*)\|_2$.}}
%	\label{error1}
%\end{figure}
%Starting from $x_0=(2,2,2)$, the plots in Figure \ref{error2} illustrate these errors for different values of $\gamma$, while the results in Figure \ref{error3} show their evolution for a fixed $\gamma=0.3$ and various initial points. 
%
%\begin{figure}[h!]
%	\subfloat{\includegraphics[scale=0.35]{diff gamma for x.png}}
%	\subfloat{\includegraphics[scale=0.35]{diff gamma for v(x).png}}
%	% Un titre pour les deux figures
%	\caption{\centering{\small The convergence rate of $\|x_k-x^*\|_2$ and $\|v(x_k)-v(x^*)\|_2$ for various $\gamma$.}}
%	\label{error2}
%\end{figure}
%
%%\medskip
%
%\begin{figure}[h!]
%	\subfloat{\includegraphics[scale=0.35]{diff x0 for x.png}}
%	\subfloat{\includegraphics[scale=0.35]{diff x0 for v(x).png}}
%	% Un titre pour les deux figures
%	\caption{\centering{\small The convergence rate of $\|x_k-x^*\|_2$ and $\|v(x_k)-v(x^*)\|_2$ for various initial points.}}
%	\label{error3}
%\end{figure}
%\end{example}
%%%%%%%%%%%%%%%%%%%%%%%%%%%%%%%%%%%%%%%%%%%%%%%%%%%%%%%%%%%%%%
%
%%
%\begin{example}
%\tcb{We illustrate Algorithm 2 in the case where the strong monotonicity is lacked. Consider
%$$
%A(x)=Qx+b,\qquad B(x)=N_{ \R_+\times \R_-\times \R_+}(x),
%$$
%where
%$$
%Q =\begin{bmatrix} 3.5 & 8 & 12.5 \\ -21 & -47 & -73 \\ 12.5 & 28 & 43.5 \end{bmatrix},  \quad
%b = \begin{bmatrix} -2 \\ 0 \\ 0 \end{bmatrix}.
%$$
%We choose 
%$$ 
%v=
%\begin{bmatrix}
%1&0&0\\
%0&-1&0\\
%0&0&1
%\end{bmatrix}.
%$$
%It is readily verified that Assumption 3 is satisfied. Moreover, the pair $ (A,v)$ is monotone and $L$-Lipschitz continuous with $L=\tcr{ missing}$.} Consequently, Algorithm 2 takes the form
%\[
%\left\{
%\begin{aligned}
%y_k&=
%P_{\mathbb R_+^3}^{\,v}
%\left(x_k-\gamma v^{-1}(Qx_k+b)\right),\\
%x_{k+1}&=
%y_k+\gamma v^{-1}(Qx_k-Qy_k). 
%\end{aligned}
%\right.
%\]
%The condition in Theorem \ref{thm4} reduces to
%$
%\gamma<\sqrt{2}.
%$
%
%Starting from the initial point
%$
%x_0=\left(3,\frac12,2\right),
%$
%the sequence generated by Algorithm 2 converges to a point in the solution set
%$S=\{(1,1,t):t\ge0\}.$ The convergence behavior is illustrated in Figure \ref{error4} for several values of $\gamma$. In particular, we plot the quantities $\|v(y_k)-v(x_k)\|_2$ and $d(y_k,S)$ as functions of the iteration number.
%
%\begin{example}
%\tcb{We illustrate Algorithm~2 in a setting where strong monotonicity is
%absent. Consider
%\[
%A(x)=Qx+b,
%\qquad
%B(x)=N_{\mathbb{R}_+\times\mathbb{R}_-\times\mathbb{R}_+}(x),
%\]
%where
%\[
%Q=
%\begin{bmatrix}
%3.5 & 8 & 12.5\\
%-21 & -47 & -73\\
%12.5 & 28 & 43.5
%\end{bmatrix},
%\qquad
%b=
%\begin{bmatrix}
%-2\\
%0\\
%0
%\end{bmatrix}.
%\]
%We choose
%\[
%v=
%\begin{bmatrix}
%1&0&0\\
%0&-1&0\\
%0&0&1
%\end{bmatrix}.
%\]
%It is readily verified that Assumption~3 is satisfied. Moreover, the pair
%$(A,v)$ is monotone and $L$-Lipschitz continuous, with
%$L=\cdots \tcr{missing}$.} Consequently, Algorithm~2 takes the form
%\[
%\left\{
%\begin{aligned}
%y_k&=
%P_{\mathbb{R}_+\times\mathbb{R}_-\times\mathbb{R}_+}^{\,v}
%\left(x_k-\gamma v^{-1}(Qx_k+b)\right),\\
%x_{k+1}&=
%y_k+\gamma v^{-1}(Qx_k-Qy_k).
%\end{aligned}
%\right.
%\]
%The stepsize condition in Theorem~\ref{thm4} reduces to
%\[
%\gamma<\sqrt{2}.
%\]
%
%Starting from the initial point
%\[
%x_0=\left(3,\frac12,2\right),
%\]
%the sequence generated by Algorithm~2 converges to a point in the
%solution set
%\[
%S=\{(1,1,t):t\geq0\}.
%\]
%The convergence behavior is illustrated in Figure~\ref{error4} for
%several values of $\gamma$. In particular, we plot the quantities
%\[
%\|v(y_k)-v(x_k)\|_2
%\qquad\text{and}\qquad
%d(y_k,S)
%\]
%as functions of the iteration number.
%\end{example}
%\begin{figure}[h!]
%	\centering{\includegraphics[scale=0.35]{expl2-algo2.png}}
%	\caption{\centering {\small The iterate error $\|v(y_k)-v(x_k)\|_2$ and the distance $d(y_k,S)$.}}
%	\label{error4}
%\end{figure}
%
%The left panel shows that the quantity $\|v(y_k)-v(x_k)\|_2$ converges rapidly to $0$. The right panel displays the distance $d(y_k,S)$, which also converges to $0$. As expected, larger values of $\gamma$  result in faster convergence..
%\end{example}

%
%\begin{figure}[h!]
%	\centering{\includegraphics[scale=0.35]{expl2-algo2.png}}
%	\caption{\centering {\small The iterate error $\|v(y_k)-v(x_k)\|_2$ and the distance $d(y_k,S)$.}}
%	\label{error4}
%\end{figure}
%
%The left panel shows that the quantity $\|v(y_k)-v(x_k)\|_2$ converges rapidly to $0$. The right panel displays the distance $d(y_k,S)$, which also converges to $0$. As expected, larger values of $\gamma$  result in faster convergence.
%\end{example}
%%%%%%%%%%%%%%%%%%%%%%%%%%%%%%%%%%%%%%%%%%%%%%%%%%%%%%%%%%
\section{Numerical Examples}\label{sec4}

\begin{example}
First we consider the operators $A: \R^3\to \R^3, B: \R^3 \rightrightarrows \R^3$ defined by
$$
A(x)=Qx+b, \qquad B(x)= N_{\R^3_+}(x),
$$
where
$$
Q =\begin{bmatrix} 1 & -3 & 0 \\ 0 & 2 & 0 \\ 0 & 0 & 1 \end{bmatrix},  \quad
b = \begin{bmatrix} -1 \\ 1 \\ -1 \end{bmatrix}.
$$
\medskip
We choose 
$$
v=\begin{bmatrix} 1 & 0 & 0\\ 0 & 3 & 0 \\ 0 & 0 & 3\end{bmatrix}.
$$
Then $A$ is non-monotone,  $v$ is invertible, $(B,v)$ is monotone, $(A,v)$ is $\alpha$-strongly monotone with $\alpha=\frac{7-\sqrt{34}}{2}\approx0.58$ and $A$ is $L$-Lipschitz continuous with {$L\approx3.71.$} Moreover, both $v$ and $v^{-1}$ are Lipschitz continuous. Hence all assumptions of Theorem \ref{thm1} are satisfied, and the convergence result follows. 
A direct computation yields the unique solution $x^*= (1,0,1)$.
\medskip

\noindent

To illustrate the convergence behavior of Algorithm 1, we compute the errors $\|x_k-x^*\|_2$ and $\|v(x_k)-v(x^*)\|_2$. Starting from $x_0=(4,3,5)$ and using the admissible stepsize $\gamma=0.06$,  the corresponding convergence behavior is displayed in Figure \ref{error1}.
%the plots in Figure \ref{error1} illustrate both errors. 
\begin{figure}[H]
	\centering{\includegraphics[scale=0.35]{1}}
	\caption{\centering {\small Iterate errors $\|x_k-x^*\|_2$ and $\|v(x_k)-v(x^*)\|_2$ for $\gamma=0.06$ and $x_0=(4,3,5)$.}}
	\label{error1}
\end{figure}

Figure \ref{error2} further illustrates the effect of the stepsize and the initial point. In the left panel, starting from $x_0=(4,3,5)$, we plot the error $\|x_k-x^*\|_2$ for different values of $\gamma$, while right panel shows its evolution for $\gamma=0.06$ and various initial points. 

\begin{figure}[H]
	\centering{\includegraphics[scale=0.35]{2}}
	\caption{\centering {\small Iterate error $\|x_k-x^*\|_2$ for various values of $\gamma$ and initial points.}}
	\label{error2}
\end{figure}
\end{example} 


%%%%%%%%%%%%%%%%%%%%%%%%%%%%%%%%%%%%%%%%%%%%%%%%%%%%%%%%%%%%%%
\begin{example}
Next we consider the operators $A: \R^3\to \R^3, B: \R^3 \rightrightarrows \R^3$ defined by
$$
A(x)=Qx+b, \qquad B(x)= N_{\R_+}(x_1)\times N_{\R_+}(x_2)\times \{N_{\R_+}(x_3)+x_3+2x_1-3x_2\},
$$
where
$$
Q =\begin{bmatrix} 1 & -3 & 0 \\ 0 & 2 & 0 \\ 0 & 0 & 0 \end{bmatrix},  \quad
b = \begin{bmatrix} -1 \\ 1 \\ 0 \end{bmatrix}.
$$
We choose 
$$
v=\begin{bmatrix} 1 & 0 & 0\\ 0 & 3 & 0 \\ 0 & 0 & 0\end{bmatrix}.
$$
Then $A$ and $B$ are non-monotone but  $(A,v)$ and $(B,v)$ are monotone and 
$$
\ran  A = {\R^2\times \{0\}= \ran v}\subset \ran (\gamma B+v)=\R^3.
$$
{Moreover $(A,v)$ is $L$-Lipschitz continuous with $L\approx 1.5$.}  Moreover, $A$ is continuous and $B$ has a closed graph, which implies that $\operatorname{gra}(A+B)$ is sequentially weak-to-strong closed. In addition,
$
S=(A+B)^{-1}(0)=\{x^*\},
\;
x^*=(1,0,0),
$
and {$(A+B)^{-1}$ is $R$-Lipschitz continuous at $0$ with Lipschitz constant $1$ and radius $\sigma=1/2$.}
%Indeed, let
%	$y=(y_1,y_2,y_3)\in\mathbb{R}^3$ with $\|y\|<\frac12$ and let
%	$x\in(A+B)^{-1}(y)$. Then
%	\[
%	\begin{cases}
%	y_1\in x_1-3x_2-1+N_{\mathbb{R}_+}(x_1),\\
%	y_2\in 2x_2+1+N_{\mathbb{R}_+}(x_2),\\
%	y_3\in N_{\mathbb{R}_+}(x_3)+x_3+2x_1-3x_2.
%	\end{cases}
%	\]
%	Since $|y_2|<\frac12$, we necessarily have $x_2=0$. The first
%	inclusion then yields $x_1=1+y_1>0$, while the third one implies
%	$x_3=0$. Hence
%	\[
%	(A+B)^{-1}(y)
%	= \left\{(1+y_1,0,0)\right\},
%	\qquad
%	\|y\|<\frac12.
%	\]
%	Consequently,
%	\[
%	\operatorname{dist}\bigl((A+B)^{-1}(y),(A+B)^{-1}(0)\bigr)
%	=|y_1|\leq\|y\|,
%	\]
%	so $(A+B)^{-1}$ is $R$-continuous at $0$, with modulus
%	$\rho(r)=r$ and radius $\sigma=\frac12$.
%with the explicit continuity modulus
%$
%\rho(r)=r
%$
%and radius
%$
%\sigma=\frac{1}{2}.
%$
Thus all assumptions of Theorem \ref{thm4} are satisfied.

Starting from the initial point $x_0=(2,1,1),$ the sequence generated by Algorithm 2 converges to $x^*$. Figure \ref{error3} illustrate this convergence for $\gamma=0.1,0.2,0.4,$ and $0.6$, all satisfying $\gamma<1/L$. 

\begin{figure}[H]
	\centering{\includegraphics[scale=0.35]{3.png}}
	\caption{\centering {\small Iterate errors $\|v(y_k)-v(x_k)\|_2$, $d(y_k,S)$,  $\|x_k-x^*\|_2$ for different values of $\gamma$ and initial points.}}
	\label{error3}
\end{figure}

The panels show that the quantities $\|v(y_k)-v(x_k)\|_2$, $d(y_k,S)$, and $\|x_k-x^*\|_2$ all converge to zero, confirming numerically the convergence established by Theorem \ref{thm4}. Moreover, larger admissible values of $\gamma$ lead to faster numerical convergence.
\end{example}

%%%%%%%%%%%%%%%%%%%%%%%%%%%%%%%%%%%%%%%%%%%%%%%%%%%%%%%%%%%%%

%Indeed, for $h=x-y$,
%\[
%\langle A(x)-A(y),v(x-y)\rangle
%=h_1^2-3h_1h_2+6h_2^2+3h_3^2,
%\]
%while
%\[
%\|v(x-y)\|_2^2=h_1^2+9h_2^2+9h_3^2.
%\]
%Consequently,
%\[
%\langle A(x)-A(y),v(x-y)\rangle
%\geq
%\frac{5-\sqrt{10}}{6}\|v(x-y)\|_2^2.
%\]
%Thus $(A,v)$ is strongly monotone with
%\[
%\alpha=\frac{5-\sqrt{10}}{6}\approx0.306287.
%\]
%Furthermore, $A$ is Lipschitz continuous with
%\[
%L=\|Q\|_2\approx3.702459.
%\]
%Hence
%\[
%\frac{2\alpha}{L^2}\approx0.044687.
%\]
%In particular, the choice $\gamma=0.04$ satisfies
%$
%0<\gamma<\frac{2\alpha}{L^2}.
%$
%
%Therefore, all the assumptions of Theorem \ref{thm1} are fulfilled
%for this choice of $\gamma$, and the sequence generated by
%Algorithm 1 converges to the unique solution.
%
%A direct computation gives
%\[
%x^*=(1,0,1)^T.
%\]
%Indeed,
%\[
%A(x^*)=(0,1,0)^T
%\quad\text{and}\quad
%-A(x^*)=(0,-1,0)^T
%\in N_{\mathbb{R}^3_+}(x^*).
%\]
%
%To illustrate the convergence behavior of Algorithm 1, we consider
%the errors
%\[
%\|x_k-x^*\|_2
%\qquad\text{and}\qquad
%\|v(x_k)-v(x^*)\|_2.
%\]
%Starting from $x_0=(4,3,5)^T$ and using the admissible stepsize
%$\gamma=0.04$, the corresponding convergence behavior is displayed
%in Figure \ref{error1}.
%%%%%%%%%%%%%%%%%%%%%%%%%%%%%%%%%%%%%%%%%%%%%%%%%%%%%%%%%%%%

%%%%%%%%%%%%%%%%%%%%%%%%%%%%%%%%%%%%%%%%%%%%%%%%%%%%%%%%%%%%%
\section{Conclusion}\label{sec5}
{This work extends the recently introduced framework of pair monotonicity from nonmonotone single-operator inclusion problems to structured sum inclusion problems involving a single-valued operator and a set-valued operator. It demonstrates that pair monotonicity provides a suitable framework for the development of generalized  splitting methods beyond the classical monotone setting. 
Two algorithms based on the transformed and warped resolvents were proposed and analyzed. Weak, strong and linear convergence results were established under suitable assumptions. Numerical examples confirmed the theoretical convergence properties and illustrated the behavior of the proposed methods.
The present work opens several perspectives for future research. These include inertial and accelerated variants and applications to non-convex constrained optimization.}

\begin{thebibliography}{99}
\bibitem{acl}
 {\sc S.~Adly, M.~G. Cojocaru, B.~K. Le},
\emph{State-dependent sweeping processes: asymptotic behavior and algorithmic approaches},
 { J. Optim. Theory Appl.}, 202(2):932--948, 2024.

\bibitem{Apidopoulos}   {\sc V. Apidopoulos, J. Aujol, C. Dossal}, \emph{Convergence rate of inertial Forward-Backward algorithm beyond Nesterov's rule}, Math Program 180, 137--156, 2020. 

\bibitem{Attouch}   {\sc  H. Attouch, A. Cabot}, \emph{Convergence Rates of Inertial Forward-Backward Algorithms}, SIAM J Optim, 28(1), 849--874, 2018.

\bibitem{Bauschke}   { \sc H. H. Bauschke, P. L. Combettes}. \emph{Convex Analysis and Monotone Operator
Theory in Hilbert Spaces}. Springer Berlin, 2011.

\bibitem{Bauschke1}   {\sc H. H. Bauschke, W. M. Moursi}.   \emph{On the Douglas-Rachford algorithm.}  Math Program, 164--263, 2017.

\bibitem{bc}
 {\sc M.~N. Bui, P.~L. Combettes},
 \emph{Warped proximal iterations for monotone inclusions}.
 { J. Math. Anal. Appl.}, 491(1):124315, 21, 2020.

\bibitem{Chen}   {\sc G. H. G. Chen, R. T. Rockafellar}, \emph{Convergence Rates in Forward-Backward Splitting}, SIAM J Optim, 7(2), 421--444, 1997.

\bibitem{cp}
 {\sc P. L. Combettes, J. Pesquet}, \emph{A Douglas-Rachford splitting approach to non-
smooth convex variational signal recovery}. IEEE J. Selected Topics in Signal Processing, 1(4):564--574, 2007.

\bibitem{cw}
 {\sc P. L. Combettes, V. R. Wajs}, \emph{Signal recovery by proximal forward-backward
splitting}, SIAM Multiscale Modeling and Simulation, 4(4):1168, 2005.




\bibitem{Davis}   {\sc D. Davis, W. Yin},  \emph{ Faster Convergence Rates of Relaxed Peaceman-Rachford and ADMM Under Regularity Assumptions}, Math Oper Res 42(3), 577--896, 2017.

\bibitem{Gibali}   {\sc A. Gibali, D. V. Thong},  \emph{Tseng type methods for solving inclusion problems and its applications}, Calcolo 55(4), 55--49, 2018.

\bibitem{Giselsson}   {\sc P. Giselsson},  \emph{Tight global linear convergence rate bounds for Douglas-Rachford splitting},
J Fixed Point Theory Appl  19,  2241--2270, 2017. 


\bibitem{Giselsson1}   {\sc P. Giselsson, S. Boyd},  \emph{Linear convergence and metric selection for Douglas-
Rachford splitting and ADMM,} IEEE Trans Automat Contr 62(2), 532-544,
2017.

\bibitem{He}   {\sc B. He, X. Yuan}, \emph{On the $o(1/n)$ convergence rate of the Douglas-Rachford
alternating direction method}, SIAM J Numer Anal 50(2), 700--709, 2012.

\bibitem{Hong}   {\sc M. Hong, Z. Q. Luo},  \emph{On the linear convergence of the alternating direction
method of multipliers}, Math Programm 162(1), 165--199, 2016.

\bibitem{Le}
{\sc  B.~K. Le},
 \emph{${R}$-continuity with applications to convergence analysis of
  {T}ikhonov regularization and {DC} programming.}
 { J. Convex Anal.}, 31(1):243--254, 2024.




\bibitem{LDT}
{\sc B.~K. Le, Minh N. Dao,  M. Théra},
\emph{Solving Non-Monotone Inclusions Using Monotonicity of Pairs of Operators},
 { J. Optim. Theory Appl.}, 207(18), 2025.


	\bibitem{LMT}
	{\sc B.~K. Le, Z. Mazgouri,  M. Théra},
	\emph{Monotonicity of Pairs of Operators and  Generalized Inertial Proximal Method},
	 {arXiv:2601.12738, 2026.} 

\bibitem{LMT1}
{\sc B. K. Le, B. S. Mordukhovich, and M. Th\'era},
\emph{Compact R-continuity with applications to solving inclusions and convergence of algorithms}, {
Set-Valued Var. Anal.},
{ 34}:13 (2026).
	 
\bibitem{LT}
{\sc B.~K. Le, M.~Th\'era},
\emph{Explicit {Convergence} {Rate} of {The} {Proximal} {Point} {Algorithm}
  under {R}-{Continuity}.}
 { Evol. Equ. Control Theory}, 17:93--105, 2026.

\bibitem{Lions}   {\sc  P. L. Lions, B. Mercier},  \emph{Splitting algorithms for the sum of two nonlinear operators,} SIAM J Numer
Anal 16, 964--979, 1979.

\bibitem{Moursi} {\sc  W. M. Moursi,  L. Vandenberghe},   \emph{Douglas-Rachford Splitting for the Sum of a Lipschitz Continuous and a Strongly Monotone Operator}, J Optim Theory Appl 183, 179--198, 2019. 



%\bibitem{Nesterov3}
%{\sc Y.~Nesterov},
% {\em Lectures on convex optimization}, volume 137 of {\em Springer
%  Optimization and Its Applications}.
% Springer, Cham, second edition, 2018.

%\bibitem{Nesterov1}
%Y.~Nesterov and B.~T. Polyak.
%\newblock Cubic regularization of {N}ewton method and its global performance.
%\newblock {\em Math. Program.}, 108(1):177--205, 2006.
%
%\bibitem{Carmon}
%L.~Nu\~nez, R.G. Regis, and K.~Varela.
%\newblock Accelerated random search for constrained global optimization
%  assisted by radial basis function surrogates.
%\newblock {\em J. Comput. Appl. Math.}, 340:276--295, 2018.

\bibitem{Opial}
 {\sc Z.~Opial},
\emph{  Weak convergence of the sequence of successive approximations for
  nonexpansive mappings},
 { Bull. Amer. Math. Soc.}, 73:591--597, 1967.

\bibitem{Svaiter}   {\sc B. F. Svaiter},\emph{ On weak convergence of the Douglas-Rachford method}, SIAM J Control Optim 49, 280--287, 2011.

\bibitem{Tseng}   {\sc P. Tseng}, \emph{A modified forward-backward splitting method for maximal monotone mappings}, SIAM J Control Optim 38(2), 431--446, 2006.


%\bibitem{Pennanen}
%T.~Pennanen.
%\newblock Local convergence of the proximal point algorithm and multiplier
%  methods without monotonicity.
%\newblock {\em Math. Oper. Res.}, 27(1):170--191, 2002.

%\bibitem{Polyak}
%B.~T. Polyak.
%\newblock Gradient methods for minimizing functionals.
%\newblock {\em \v Z. Vy\v cisl. Mat i Mat. Fiz.}, 3:643--653, 1963.
%
%\bibitem{Rockafellar}
%R.~T. Rockafellar.
%\newblock Monotone operators and the proximal point algorithm.
%\newblock {\em SIAM J. Control Optim.}, 14(5):877--898, 1976.
%
%\bibitem{Rockafellar2019}
%R.~T. Rockafellar.
%\newblock Progressive decoupling of linkages in optimization and variational
%  inequalities with elicitable convexity or monotonicity.
%\newblock {\em Set-Valued Var. Anal.}, 27(4):863--893, 2019.
%
%\bibitem{Pham}
%P.~D. Tao and L.~T.~H. An.
%\newblock Convex analysis approach to d.c. programming: theory, algorithms and
%  applications.
%\newblock {\em Acta Math. Vietnam.}, 22(1):289--355, 1997.

%\bibitem{Tikhonov}
%A.~Tikhonov and V.~Ars\'enine.
%\newblock {\em M\'ethodes de resolution de probl\`emes mal pos\'es}.
%\newblock \'Editions Mir, Moscow, 1976.
%\newblock Traduit du russe par Vladimir Kotliar.

%\bibitem{Wittmann}
%R.~Wittmann.
%\newblock Approximation of fixed points of nonexpansive mappings.
%\newblock {\em Arch. Math. (Basel)}, 58(5):486--491, 1992.

\end{thebibliography}

\end{document}



%%%%%%%%%%%%%%%%%%%%%%%%%%%%%%%%%%%%%%%%%%%%%%%%%%%%%%%%%%%%%%%%%%%%%%%%%%%%%%%%%%%%%%%%%%%%%%%%%%



\end{document}

%%%%%%%%%%%%%%%%%%%%%%%%%%%%%%%%%%%%%%%%%%%%%%%%%%%%%%%%%%%%%
\subsection*{Example 2}

%%%%%%%%%%%%%%%%%%%%%%%%%%%%%%%%%%%%%

We illustrate Algorithm 2 in a case where \(A\) is not surjective. Consider
\[
A(x)=Qx+b,\qquad
B(x)=N_{\mathbb{R}_+^3}(x),
\]
where
\[
Q=
\begin{pmatrix}
1&1&0\\
-1&1&0\\
0&0&0
\end{pmatrix},
\qquad
b=
\begin{pmatrix}
-2\\0\\0
\end{pmatrix},
\qquad
v(x)=Vx,\ \
V=
\begin{pmatrix}
2&0&0\\
0&2&0\\
0&0&1
\end{pmatrix}.
\]
Since
\[
\operatorname{rge}(A)=\mathbb{R}^2\times\{0\}
\subset
\operatorname{rge}(v)=\mathbb{R}^3,
\]
Assumption 3 is satisfied. Moreover,
\[
\langle A(x)-A(y),v(x)-v(y)\rangle
=2(z_1^2+z_2^2)\ge0,
\quad z=x-y,
\]
so that \((A,v)\) is monotone. In addition,
\[
\|QV^{-1}\|=\frac1{\sqrt2},
\]
and hence \((A,v)\) is \(L\)-Lipschitz continuous with
$
L=\frac1{\sqrt2}.
$

The solution set is
\[
S=\{(1,1,t):t\ge0\},
\]
and Algorithm 2 reduces to
\[
\left\{
\begin{aligned}
y_k&=
P_{\mathbb R_+^3}^{\,v}
\left(x_k-\gamma V^{-1}(Qx_k+b)\right),\\
x_{k+1}&=
y_k+\gamma V^{-1}(Qx_k-Qy_k).
\end{aligned}
\right.
\]
The condition of Theorem \ref{thm4} becomes
\[
0<\gamma<\sqrt2,
\]
and the numerical experiments are performed for admissible values of \(\gamma\). The convergence is assessed through 
\[
\|v(y_k)-v(x_k)\|_2
\]
and the distance to the solution set
\[
d(y_k,S)=\sqrt{(y_1^k-1)^2+(y_2^k-1)^2}.
\]

Starting from 
$
x_0=(3,\,0.5,\,2),
$
the sequence \((x_k)\) generated by Algorithm 2 converges to a point of the solution set \(S\).
%%%%%%%%%%%%%%%%%%%%%%%%%%%%%%%%%%%%%%%%%%%%%%%%%%%%%%%%%%%%
%%%%%%%%%%%%%%%%%%%%%%%%%%%%%%%%%%%%%%%%%%%%%%%%%%%%%%%%%
%%%%%%%%%%%%%%%%%%%%%%%%%%%%%%%%%%%%%%%%%%%%%%%%%%%%%%%%%%%%%%%%%%%%





